\documentclass[10pt,a4paper]{article}
\usepackage[margin=0.5in]{geometry}
\usepackage[T1]{fontenc}
\usepackage{lmodern}
\usepackage{booktabs}
\usepackage{tabularx}
\usepackage{longtable}
\usepackage{array}
\usepackage{enumitem}
\usepackage{microtype}
\usepackage[hidelinks]{hyperref}
\setlist{itemsep=2pt,topsep=3pt}
\pagestyle{empty}
\setlength{\parindent}{0pt}
\setlength{\parskip}{4pt}
\newcolumntype{Y}{>{\raggedright\arraybackslash}X}
\renewcommand{\arraystretch}{1.15}
\newcommand{\code}[1]{\texttt{#1}}

\title{Baseline Retrieval vs.\ Cross-Encoder Re-Ranking\\[2pt]\large A RAG comparison on the ICC Playing Handbook 2019--20}
\author{}
\date{}

\begin{document}
\maketitle
\thispagestyle{empty}
\vspace{-3em}

\section*{Setup}

\textbf{Document}
\begin{itemize}
\item ICC Playing Handbook 2019--20 (cricket rules), 427 PDF pages.
\item About 297,000 words after cleaning, split into 2,660 chunks of at most 150 words.
\item The format name (for example ``ICC Men's One Day International Playing Conditions'') is written at the start of each chunk.
\end{itemize}

\textbf{Models}
\begin{itemize}
\item Bi-encoder: \code{all-MiniLM-L6-v2} (384-dimensional vectors, FAISS flat index).
\item Cross-encoder: \code{ms-marco-MiniLM-L-6-v2}.
\item LLM: \code{llama3.1:8b} through Ollama, temperature 0.
\end{itemize}

\textbf{Pipelines}
\begin{itemize}
\item \textbf{Pipeline 1 (baseline):} bi-encoder top 10, then the top 3 go to the LLM $\rightarrow$ Answer 1.
\item \textbf{Pipeline 2 (re-ranked):} the same top 10 are re-scored by the cross-encoder, sorted, and the top 3 go to the LLM $\rightarrow$ Answer 2.
\item The prompt is identical in both pipelines.
\end{itemize}

\textbf{Machine}
\begin{itemize}
\item MacBook Air M4, 24 GB memory; everything runs locally.
\end{itemize}

\section{Architecture: bi-encoder vs.\ cross-encoder}

\subsection*{Bi-encoder}
\textbf{How it works}
\begin{itemize}
\item The same transformer reads the query and each chunk \emph{separately}.
\item Each text becomes one vector of 384 numbers (the average of the token vectors, scaled to length 1).
\item The score is the dot product of the two vectors, which equals the cosine similarity.
\end{itemize}
\textbf{Speed}
\begin{itemize}
\item Chunk vectors are computed once and stored.
\item A search needs one pass for the query, then fast dot products.
\end{itemize}
\textbf{Limits}
\begin{itemize}
\item A chunk is squeezed into 384 numbers \emph{before the question is known}, so small details (20 vs.\ 50, ``not'') can be blurred.
\item Input limit: 256 tokens.
\end{itemize}
\newpage
\subsection*{Cross-encoder}
\textbf{How it works}
\begin{itemize}
\item The query and one chunk are joined into \emph{one input}: \code{[CLS] query [SEP] chunk [SEP]}.
\item In every layer, self-attention lets each query token look at each chunk token and the other way round.
\item The \code{[CLS]} output goes to a small layer that gives \emph{one relevance score} for the pair. It is not a cosine, so scores can only be compared within one query.
\end{itemize}
\textbf{Speed}
\begin{itemize}
\item Nothing can be computed in advance: each (query, chunk) pair needs a full pass.
\end{itemize}
\textbf{Limits}
\begin{itemize}
\item Input limit: 512 tokens for the pair.
\item More accurate on details, but slower.
\end{itemize}

\subsection*{Side by side}
\begin{center}\small
\begin{tabular}{@{}l l l@{}}
\toprule
 & Bi-encoder & Cross-encoder \\
\midrule
Input & Query and chunk, separately & Query and chunk, together \\
Output & One vector per text & One score per pair \\
Score & Cosine (dot product) & Raw relevance score \\
Chunk work done in advance & Yes (vectors stored) & No \\
Cost per query & 1 pass + dot products & 1 pass per chunk \\
Input limit & 256 tokens & 512 tokens (pair) \\
Used here for & Top 10 from 2,660 chunks & Re-sorting those 10 \\
\bottomrule
\end{tabular}
\end{center}

\section{Rank shifts and answer comparison}

Rank is the position of the right chunk in the top 10 (bi-encoder $\rightarrow$ cross-encoder). Chunk id and page are in brackets. ``Omitted'' means the LLM replied ``The provided excerpts do not contain the answer.''

{\small\renewcommand{\arraystretch}{1.0}
\begin{longtable}{@{}r >{\raggedright}p{1.8cm} >{\raggedright}p{4.2cm} >{\raggedright}p{3.0cm} >{\raggedright}p{2.0cm} >{\raggedright\arraybackslash}p{3.2cm}@{}}
\toprule
\# & Type & Query (short) & Right chunk, rank & Answer 1 & Answer 2 \\
\midrule
\endhead
1 & Buried chunk & Overs per team, women's T20I & \#1626 (p9.26): $7\rightarrow1$ & Omitted & ``20 overs''. Correct \\
2 & Buried chunk & Fielders outside the circle, overs 11--40, men's ODI & \#589 (p4.55): $6\rightarrow1$ & Omitted & ``Four''. Correct \\
3 & Buried chunk & Overs per innings, men's ODI & \#473 (p4.24): $10\rightarrow2$ & Omitted & ``50 overs''. Correct \\
4 & Exception & Dismissal on a free hit, men's T20I & \#924 (p5.43): $7\rightarrow1$ & Omitted & Only in the ways allowed for a No ball. Correct \\
5 & Negation & Can a suspended bowler bowl again in the innings, men's ODI & \#646 (p4.73): $8\rightarrow1$ & Omitted & ``No, shall not be allowed to bowl again in that innings''. Correct \\
6 & Keyword / synonym & ``Short balls at head height'' per over, men's T20I (rule says ``fast short-pitched delivery'') & \#1017 (p5.70): $7\rightarrow1$ & Omitted & Omitted (chunk found, LLM still refused) \\
7 & Re-ranking hurts & Fielders outside the circle, first ten overs, men's ODI & \#588 (p4.55): $1\rightarrow8$ & ``Two''. Correct & Omitted \\
8 & Re-ranking hurts & Bat straight after 8+ min off the field, men's T20I & \#937 (p5.47): $3\rightarrow10$ & ``No, not until\dots''. Correct & Omitted \\
9 & Not in top 10 & Overs per bowler, men's T20I & \#862 (p5.26): not in top 10 & Omitted & Omitted \\
10 & Not in top 10 & Runner for an injured batter, T20I & \#944 (p5.48): not in top 10 & Omitted & Omitted \\
\bottomrule
\end{longtable}}

\newpage
\textbf{What the results show.}
\begin{itemize}
\item \textbf{Why the baseline fails.} The five formats (men's Test, ODI, T20I, women's Championship, women's T20I) repeat most clauses with small changes, for example 50 or 20 overs. The baseline top 3 often held chunks from the wrong format or the wrong Powerplay phase. In query 2 it held the first-Powerplay ODI chunk (\#588), a T20I chunk (\#964) and a women's Championship chunk (\#1356), but not the right one.
\item \textbf{The baseline omitted the clause; it did not make up a number.} The prompt tells the LLM to use only the excerpts, so a missed chunk became a refusal. In queries 1--5 this is the ``omits an essential clause'' case. No hallucination was seen in these 10 runs.
\item \textbf{The cross-encoder fixed 6 of 10} (queries 1--6). It reads the question together with the format name in the chunk, so it can pick the right copy. In queries 1--5 this also gave a correct Answer 2.
\item \textbf{It made 2 worse} (queries 7, 8): it ranked a chunk for another phase higher, and the right chunk left the top 3.
\item \textbf{It cannot help when the chunk is not in the top 10} (queries 9, 10). The rule is one sentence inside a long chunk, or the question uses a word the document does not (for example ``bouncers''). The cross-encoder only re-sorts the top 10.
\item \textbf{Query 6:} retrieval was fixed but the 8B LLM still refused, so not every failure is a retrieval failure.
\end{itemize}

\textbf{Limit of this test.} About 80 questions were tried and these 10 were kept. Most questions were already answered correctly by both pipelines. The results show \emph{when} re-ranking helps or hurts, not how often.

\section{Latency profiling}

Measured on a MacBook Air M4 (24 GB), 10 queries $\times$ 10 repeats, after a warm-up. Mean $\pm$ standard deviation, in milliseconds. The top 10 comes from an exact (flat) FAISS index over 2,660 chunks.

\begin{center}\small
\begin{tabular}{@{}l r r@{}}
\toprule
Stage & Pipeline 1 (ms) & Pipeline 2 (ms) \\
\midrule
Stage 1: embed query + FAISS top 10 & $5.0 \pm 1.0$ & $5.0 \pm 1.0$ \\
Stage 2: cross-encoder on 10 pairs & not used & $35.9 \pm 3.3$ \\
\textbf{Retrieval total} & $\mathbf{5.0 \pm 1.0}$ & $\mathbf{40.9 \pm 3.8}$ \\
LLM answer (\code{llama3.1:8b}, top 3 chunks) & $1379 \pm 1173$ & $1700 \pm 1254$ \\
\bottomrule
\end{tabular}
\end{center}

Re-ranking adds about 36 ms, about 8 times the retrieval time. This is small next to the LLM. LLM time varies a lot because it depends on answer length (a one-line refusal takes about 0.5 s).

\section{Why not run the cross-encoder on the whole collection?}

\begin{enumerate}
\item \textbf{Nothing can be precomputed.} A bi-encoder embeds each chunk once, before any query. A cross-encoder needs the query, so every (query, chunk) pair needs a new transformer pass for every query.
\item \textbf{Cost grows with collection size.} Measured: 35.9 ms for 10 pairs, about 3.6 ms per pair. For our 2,660 chunks that is about 9.5 s per query. For 1 million chunks it would be about 1 hour per query. (This is an estimate from the measured time; a GPU would be faster, but the growth is the same.) The bi-encoder search over 2,660 chunks takes 5 ms, and large systems use approximate indexes that look at only a small part of the vectors.
\item \textbf{So we use two stages.} The cheap bi-encoder picks a small candidate set (top 10); the expensive cross-encoder re-orders only those. The limit: a chunk the bi-encoder did not put in the top 10 can never be recovered (queries 9, 10). A larger top-$K$ helps, but cost grows with $K$.
\end{enumerate}

\end{document}
